% Propietario conservado: /Users/ruben/Documents/ChatGPT/jueces y controles/REVISION_ARGUMENTAL_APERTURAS_CIERRES_20260915/delivery/ARTICLE_V_ES_ARGUMENTAL_20260915/manuscrito/sections/61_incidencia_entropia.tex
\section{Incidence surfacique--volumétrique et entropie topologique}

Le calendrier du TPK induit la charnière entre les modes additif et
multiplicatif. Dans le bloc primitif
\[
 \tau_3=(+1,-1,0),
\]
la règle
\[
 +1:m\mapsto\Sigma,
 \qquad
 -1:m\mapsto\Pi,
 \qquad
 0:m\mapsto m
\]
produit l’orbite cyclique \((\Sigma,\Pi,\Pi)\). Le lecteur géométrique
\[
 \Sigma\longmapsto\mathscr H_{\rm ar},
 \qquad
 \Pi\longmapsto\mathscr H_{\rm vol}
\]
transporte ces modes vers les feuillets surfacique et volumétrique.

\begin{theorem}[Incidence nécessaire des feuillets]
\label{v:thm:incidencia-hojas-anexo}
Dans l’ordre \((\mathscr H_{\rm ar},\mathscr H_{\rm vol})\), la matrice
primitive des incidences est
\begin{equation}
 \boxed{
 F_{\rm av}=
 \begin{pmatrix}0&1\\1&1\end{pmatrix}.}
 \label{v:eq:fav-anexo-informacion}
\end{equation}
Les calendriers de longueurs \(9,27,108\) produisent
\[
 N_9=3F_{\rm av},
 \qquad N_{27}=9F_{\rm av},
 \qquad N_{108}=36F_{\rm av}.
\]
\end{theorem}

\begin{proof}
Les paires consécutives du cycle sont
\(\Sigma\to\Pi\), \(\Pi\to\Pi\) et \(\Pi\to\Sigma\). Chacune apparaît une
fois et \(\Sigma\to\Sigma\) n’apparaît pas. Cela fixe les quatre entrées de
\(F_{\rm av}\). Les calendriers longs contiennent, respectivement, trois,
neuf et trente-six copies du bloc primitif.
\end{proof}

\begin{theorem}[Loi de Fibonacci et mesure d’entropie maximale]
\label{v:thm:fibonacci-parry-anexo}
Pour \(n\geq1\),
\begin{equation}
 F_{\rm av}^{n}
 =\begin{pmatrix}F_{n-1}&F_n\\F_n&F_{n+1}\end{pmatrix}.
 \label{v:eq:potencias-fav-anexo}
\end{equation}
Le sous-décalage à mémoire un défini par \(F_{\rm av}\) a
\begin{equation}
 \boxed{h_{\rm top}=\log\varphi},
 \qquad
 \boxed{\frac{\mu_{\rm ar}}{\mu_{\rm vol}}=\varphi^{-2}}.
 \label{v:eq:entropia-parry-hojas-anexo}
\end{equation}
\end{theorem}

\begin{proof}
L’identité matricielle s’obtient par récurrence à partir de
\(F_{n+1}=F_n+F_{n-1}\). La racine de Perron de \(F_{\rm av}\) est \(\varphi\),
d’où \(h_{\rm top}=\log\varphi\). Comme la matrice est symétrique, les poids
de Parry \cite{v:Parry1964} sont proportionnels au carré d’un vecteur de Perron
\((1,\varphi)\), c’est-à-dire à \((1,\varphi^2)\).
\end{proof}

Les expressions \(\log3\) et \(\log\varphi\) sont ainsi séparées par leur
généalogie. La première quantifie un choix local uniforme entre trois
états ; la seconde mesure la croissance des histoires permises par
l’incidence des feuillets. La relation surfacique--volumétrique est une propriété du
retour complet et apparaît donc avec l’exposant deux.

